\documentclass[11pt]{article}
\usepackage[T1]{fontenc}
\usepackage{textcomp}
\usepackage[margin=0.75in]{geometry}
\usepackage{amsmath,amssymb,amsthm}
\usepackage{array,booktabs}
\usepackage{hyperref}
\setlength{\parindent}{0pt}
\newtheorem{theorem}{Theorem}
\theoremstyle{definition}
\newtheorem{definition}{Definition}
\title{Flow Proof Artifact: Monoid-identity-unique}
\author{Generated by \texttt{flow doc proof}}
\date{Flow Proof Book}
\begin{document}
\maketitle
The identity element of a monoid is unique.\\[0.5em]
\textbf{Source.} dummit-foote — *Abstract Algebra*, §1.1\\[0.5em]
\bigskip
\noindent{\large \textbf{Derived fact 1.} \textit{If e and e' are both identities then e = e'} \hfill Monoid \cdot identity \cdot the identity is unique}\par
\smallskip\noindent\textit{Coordinate: Monoid \cdot identity \cdot the identity is unique}\par
\smallskip\noindent\textit{Source: dummit-foote}\par
\smallskip\noindent\textit{Built on: one is the left identity, for identity on Monoid, one is the right identity, for identity on Monoid, parentheses do not matter, for multiplication on Monoid}\par
\medskip
\noindent\textbf{Goal.} If e and e' are both identities then e = e'\par
\noindent$\displaystyle \forall e \in Monoid \forall e_prime \in Monoid\quad e = \text{e prime}$\par
\medskip
\noindent
\begin{tabular}{@{} >{\raggedright\arraybackslash}p{0.06\textwidth} >{\raggedright\arraybackslash}p{0.47\textwidth} >{\raggedleft\arraybackslash}p{0.41\textwidth} @{}}
\textbf{\#} & \textbf{Proof} & \textbf{Mathematics} \\
\midrule
\textbf{1.} & We prove that the identity is unique for identity on Monoid. & \\[0.45em]
\textbf{2.} & We split into exhaustive cases — the claim must hold in each one. & \\[0.45em]
\textbf{3.} & Case 1 (see step 2): suppose e * e\_prime  equals  e\_prime. & \\[0.45em]
\textbf{4.} & Case 2 (see step 2): suppose e\_prime * e  equals  e. & \\[0.45em]
\textbf{5.} & We invoke the definitional clause governing identity on Monoid: one is the left identity, for identity on Monoid (instantiated for e\_prime). & \\[0.45em]
\textbf{6.} & We invoke the definitional clause governing identity on Monoid: one is the right identity, for identity on Monoid (instantiated for e). & \\[0.45em]
\textbf{7.} & We invoke the definitional clause governing multiplication on Monoid: parentheses do not matter, for multiplication on Monoid (instantiated for e, e\_prime, e). & \\[0.45em]
\textbf{8.} & From step 4, step 5, step 6, and step 7, this implies e equals e prime. Together with the other cases (step 3 and step 4), the goal is discharged. Hence proven. & $\displaystyle e = \text{e prime}$ \\[0.45em]
\bottomrule
\end{tabular}
\label{thm:Monoid-identity-the_identity_is_unique}
\par\medskip\hrule
\end{document}
